\documentclass[11pt]{article}
\usepackage[margin=1in]{geometry}
\usepackage{amsmath,amssymb,amsthm,booktabs,microtype}
\usepackage[most]{tcolorbox}
\usepackage[hidelinks]{hyperref}
\newtheorem{proposition}{Proposition}
\newtheorem{lemma}{Lemma}
\newtheorem{definition}{Definition}
\newtheorem{remark}{Remark}
\newtcolorbox{claimbox}[1]{colback=gray!4,colframe=black!65,title=\textbf{#1},fonttitle=\bfseries,boxrule=0.6pt,arc=1pt}
\title{\textbf{A Constructible Tower Approaching $\pi$}\\\large Explicit Compass-and-Straightedge Approximants under Transparent Algebraic Acceleration}
\author{C. Wayne Baker}
\date{Original: May 24, 2026\\Review revision v1.2: October 24, 2026}
\begin{document}
\maketitle
\begin{abstract}
We exhibit an explicit sequence of compass-and-straightedge constructible real numbers approaching $\pi_g=\pi$, although the limit itself is not constructible. The Archimedean half-perimeter $P_N=N\sin(\pi/N)$ is sampled on $N_i=4\cdot2^i$, and its values are combined using classical Richardson--Romberg extrapolation
\[
T_{i,0}=P_{N_i},\qquad T_{i,j}=T_{i,j-1}+\frac{T_{i,j-1}-T_{i-1,j-1}}{4^j-1}.
\]
Every finite tableau entry remains constructible. The diagonal has leading error
\[
\pi-T_{i,i}=\frac{\pi^{2i+3}}{(2i+3)!}\,4^{-(i+1)(i+4)/2}\bigl(1+O(i^{-2})\bigr)
\]
and correct-digit count
\[
D(i)=(\log_{10}2)i^2+2i\log_{10}i+c_1i+\tfrac72\log_{10}i+O(1),\quad
c_1=7\log_{10}2-2\log_{10}(\pi e).
\]
The increasing formal cancellation order does not imply an arithmetic speed record: the full triangular tableau uses $i(i+1)/2$ rational extrapolation updates, in addition to $i$ geometric bisection levels. The contribution is the transparent constructible provenance, interpolation error law, and separation between finite algebraic approximants and a transcendental limit.
\end{abstract}
\section{The two facts in tension}
For a regular $N$-gon inscribed in the unit circle, the half-perimeter $P_N=N\sin(\pi/N)$ is constructible when $N$ is a power of two. Set $N_i=4\cdot2^i$, $i\ge0$. Repeated bisection constructs each base length without using the numerical value of $\pi$. The limit is $\pi_g:=\lim_{i\to\infty}P_{N_i}=\pi$. Lindemann's theorem implies that $\pi$ is transcendental and therefore not compass-and-straightedge constructible. Approaching a nonconstructible limit through constructible finite objects is classical; our narrower interest is the explicit effect of rational extrapolation on this ladder.
\section{The acceleration preserves constructibility}
Define $T_{i,0}=P_{N_i}$ and, for $1\le j\le i$,
\[
T_{i,j}=T_{i,j-1}+\frac{T_{i,j-1}-T_{i-1,j-1}}{4^j-1}.
\]
The rational denominators begin $3,15,63,255,\ldots$.
\begin{proposition}[Preservation of constructibility]
Every $T_{i,j}$, in particular every diagonal $T_{i,i}$, is a constructible real number.
\end{proposition}
\begin{proof}
Each base half-perimeter $P_{N_i}$ is constructible by repeated angle bisection. The constructible real numbers form a field; the tableau uses only addition, subtraction, and division by nonzero rational integers. Induction on $j$ completes the proof.
\end{proof}
\begin{remark}[No circular use of $\pi$]
The trigonometric formula describes a length already constructed from the unit circle. It does not provide $\pi$ as an input to the construction.
\end{remark}
\section{Interpolation form and diagonal error}
Put $x_i=N_i^{-2}=4^{-(i+2)}$ and define the entire function
\[
F(x)=\begin{cases}\sin(\pi\sqrt x)/\sqrt x,&x\ne0,\\ \pi,&x=0.\end{cases}
=\sum_{m=0}^{\infty}(-1)^m\frac{\pi^{2m+1}}{(2m+1)!}x^m.
\]
Then $F(x_i)=P_{N_i}$. The recurrence is Neville--Richardson interpolation: $T_{i,i}$ is the value at $0$ of the polynomial interpolating $F$ at $x_0,\ldots,x_i$. Thus it annihilates the terms $x,\ldots,x^i$, equivalently $N^{-2},\ldots,N^{-2i}$.
\begin{lemma}[Leading diagonal error]
As $i\to\infty$,
\[
\pi-T_{i,i}=\frac{\pi^{2i+3}}{(2i+3)!}\prod_{r=0}^{i}x_r\bigl(1+O(i^{-2})\bigr)
=\frac{\pi^{2i+3}}{(2i+3)!}4^{-(i+1)(i+4)/2}\bigl(1+O(i^{-2})\bigr).
\]
In particular the diagonal eventually approaches $\pi$ from below.
\end{lemma}
\begin{proof}
The interpolant of $x^{i+1}$ evaluated at $0$ is $(-1)^i\prod_{r=0}^i x_r$. Its Taylor coefficient in $F$ is $(-1)^{i+1}\pi^{2i+3}/(2i+3)!$, giving a positive contribution to $F(0)-T_{i,i}$. For later powers the complete homogeneous symmetric polynomials are bounded by $h_k(x_0,\ldots,x_i)\le(\sum_{r=0}^ix_r)^k$, with $\sum_{r=0}^ix_r<1/12$. The absolute ratios of successive Taylor coefficients are $\pi^2/((2m+2)(2m+3))=O(i^{-2})$ for $m\ge i+1$. The normalized tail is therefore $O(i^{-2})$. Finally $\prod_{r=0}^ix_r=4^{-\sum_{r=0}^i(r+2)}=4^{-(i+1)(i+4)/2}$.
\end{proof}
\begin{definition}[Geometric circle constant]
The accelerated tower defines $\pi_g:=\lim_{i\to\infty}T_{i,i}=\pi$. The subscript records the geometric approximation pathway, not a distinct number.
\end{definition}
\section{Sharper digit asymptotics}
Write $D(i)=-\log_{10}|\pi-T_{i,i}|$. The lemma yields
\[
D(i)=\frac{(i+1)(i+4)}2\log_{10}4+\log_{10}\bigl((2i+3)!\bigr)-(2i+3)\log_{10}\pi+O(i^{-2}).
\]
\begin{proposition}[Digit-growth law with explicit constant]\label{prop:digits}
\[
\begin{aligned}
D(i)={}&(\log_{10}2)i^2+2i\log_{10}i+c_1i+\tfrac72\log_{10}i\\
&+K+\frac{73\log_{10}e}{24i}+O(i^{-2}),\\
c_1={}&7\log_{10}2-2\log_{10}(\pi e),\\
K={}&8\log_{10}2-\tfrac52\log_{10}\pi\approx1.16536528.
\end{aligned}
\]
Consequently $D(i)/i^2\to\log_{10}2$.
\end{proposition}
\begin{proof}
The geometric-node exponent is $(i+1)(i+4)=i^2+5i+4$ in base $2$. The logarithmic Stirling expansion through inverse order is
\[
\begin{aligned}
\log((2i+3)!)={}&(2i+\tfrac72)\log i+2i\log2-2i\\
&+\tfrac72\log2+\tfrac12\log(2\pi)
+\frac{73}{24i}+O(i^{-2}).
\end{aligned}
\]
Substituting this into the preceding expression for $D(i)$, the node-product contribution $(i^2+5i+4)\log_{10}2$, and $-(2i+3)\log_{10}\pi$ gives the stated linear coefficient, constant and inverse-order term. The $O(i^{-2})$ error in the leading interpolation formula contributes only $O(i^{-2})$ to $D(i)$ and does not affect these coefficients.
\end{proof}
The $2i\log_{10}i$ and lower-order terms explain why accessible ratios $D(i)/i^2$ remain above their limiting constant.
\begin{table}[ht]\centering\small
\begin{tabular}{rrrr}\toprule Level $i$ & $N_i=4\cdot2^i$ & Correct digits $D(i)$ & $D(i)/i^2$\\\midrule
5&128&19.59&0.78355\\10&4096&57.34&0.57337\\15&131072&112.05&0.49798\\20&4194304&183.12&0.45781\\25&134217728&270.26&0.43241\\30&4294967296&373.26&0.41474\\\bottomrule\end{tabular}
\caption{Original diagonal data obtained using 600-digit arithmetic; an AI-assisted 700-digit recomputation reproduced the stated rounding.}\label{tab:tower}
\end{table}
\begin{table}[ht]
\centering\small
\begin{tabular}{rrrr}
\toprule
$i$ & $D(i)$ & $R_i$ & $i^2R_i$\\
\midrule
5  & 19.5886515 & $-0.0243609$ & $-0.6090$\\
10 & 57.3372619 & $-0.00641350$ & $-0.6413$\\
15 & 112.046159 & $-0.00289698$ & $-0.6518$\\
20 & 183.123000 & $-0.00164218$ & $-0.6569$\\
25 & 270.258718 & $-0.00105571$ & $-0.6598$\\
30 & 373.262496 & $-0.000735272$ & $-0.6617$\\
\bottomrule
\end{tabular}
\caption{Numerical check of the inverse-order refinement. Here
$R_i=D(i)-[(\log_{10}2)i^2+2i\log_{10}i+c_1i+
\tfrac72\log_{10}i+K+(73\log_{10}e)/(24i)]$.
Values were recomputed using a 650-decimal-digit \texttt{mpmath} tableau.
The bounded, slowly varying $i^2R_i$ is consistent with a $O(i^{-2})$ residual;
these finite rows do not establish its limiting constant.}
\label{tab:digits-residual}
\end{table}

\section{Order, operation count, and numerical conditioning}
For each fixed $i$, the diagonal polynomial cancels through $N^{-2i}$; formally its first surviving asymptotic term has order $N^{-2i-2}$. The term ``unbounded order'' refers to this growing sequence of cancellation orders, not a fixed-cost infinite-order algorithm.
Constructing the ladder through $N_i$ requires $i$ additional bisection levels beyond the square. Populating the full lower triangular tableau requires exactly $1+2+\cdots+i=i(i+1)/2$ extrapolation updates and $i+1$ base values. These are counts of algebraic operations at growing precision, not a count of physical straightedge-and-compass moves. Since $D(i)\sim(\log_{10}2)i^2$, the asymptotic index satisfies $i\sim\sqrt{D/\log_{10}2}$ and the number of updates is $\sim D/(2\log_{10}2)$; bit complexity also depends on the working precision and the representation of constructible numbers. These observations do not establish superiority over direct polygonal doubling or modern algorithms for $\pi$.
The extrapolation uses the classical base-two Richardson weights (recurrence denominators $4^j-1$). A companion scale-optimized N-series stability study reports an aggregate absolute-weight limit $S_2\approx1.969$ for the base-two coefficients~\cite{bakerstability}; this contextualizes, but does not certify, finite-precision stability. The numerical stability of a particular finite-precision implementation must be considered separately from algebraic constructibility; a global limiting weight norm should not be interpreted as a complete floating-point error bound.
\section{AI-assistance and numerical provenance}
The author directed the study and retains responsibility for the mathematical claims. ChatGPT assisted with LaTeX revision, symbolic Stirling calculations, high-precision residual-table recomputation and manuscript preparation. Claude assisted by independently recomputing the diagonal table at 700-digit working precision and reviewing the asymptotic argument. These are AI-assisted checks, not an independent human peer review; the original table was reported from a 600-digit calculation. No archival manuscript or numerical verification results are modified by this review-copy revision.
\section{Claim boundary and conclusion}
\begin{claimbox}{Established claim}
Every finite $T_{i,i}$ is constructible and $T_{i,i}\to\pi_g=\pi$, with the explicit diagonal error law and sharpened digit-growth asymptotic with explicit constant and inverse-order correction given above.
\end{claimbox}
\begin{claimbox}{Not claimed}
The construction neither makes $\pi$ constructible nor gives a new extrapolation method or a competitive arithmetic speed record. Richardson--Romberg cancellation is classical. No historical priority is asserted for applying such tabulations to Archimedean polygon sequences.
\end{claimbox}
The tower is geometric at its base, algebraic in its acceleration, and analytic in its verification. It approaches the transcendental limit without crossing the classical constructibility boundary at any finite stage.
\newpage
\begin{thebibliography}{12}
\bibitem{bakerstability} C. W. Baker, \emph{Scale-Optimized N-Series Acceleration}, companion working manuscript on extrapolation weight norms (2026); bibliographic details to be confirmed before publication.
\bibitem{brezinski} C. Brezinski, ``Convergence acceleration during the 20th century,'' \emph{J. Comput. Appl. Math.} \textbf{122} (2000), 1--21. This historical survey documents the earlier Bauer--Stiefel--Rutishauser extrapolation literature, without establishing priority for this particular polygonal tableau.
\bibitem{archimedes} Archimedes, \emph{Measurement of a Circle}.
\bibitem{wantzel} P.-L. Wantzel, ``Recherches sur les moyens de reconna\^itre si un probl\`eme de g\'eom\'etrie peut se r\'esoudre avec la r\`egle et le compas,'' \emph{J. Math. Pures Appl.} 2 (1837), 366--372.
\bibitem{lindemann} F. Lindemann, ``\"Uber die Zahl $\pi$,'' \emph{Math. Ann.} 20 (1882), 213--225.
\bibitem{richardson} L. F. Richardson, ``The approximate arithmetical solution by finite differences of physical problems,'' \emph{Phil. Trans. Roy. Soc. A} 210 (1911), 307--357.
\bibitem{romberg} W. Romberg, ``Vereinfachte numerische Integration,'' \emph{Det Kongelige Norske Videnskabers Selskab Forhandlinger} 28 (1955), 30--36.
\end{thebibliography}
\end{document}

